\documentclass[11pt]{article}

\usepackage[margin=1in]{geometry}
\usepackage{amsmath, amssymb}
\usepackage[T1]{fontenc}
\usepackage[utf8]{inputenc}
\usepackage{lmodern}
\usepackage{enumitem}
\usepackage{hyperref}

\setlist[itemize]{topsep=4pt, itemsep=2pt, parsep=0pt}

\title{6.8610 Spring 2026 Lecture 7 Notes: Scaling Laws and Emergent Behavior}
\author{Junru Ren with help from Codex}
\date{February 26, 2026}

\begin{document}
\maketitle

\section{Scaling Laws and Emergent Behavior}

\subsection{Motivation}
This lecture studies a central practical question in modern NLP: if one has a fixed compute budget, what training run should be launched to obtain the best model? The answer motivated an entire research program around \textbf{scaling laws}, which try to predict how loss changes as model size, data size, and training compute increase.

\begin{itemize}
    \item In recent years, language modeling improved not just by better architectures, but by systematically scaling models, data, and compute.
    \item Large language models are foundation models trained with self-supervised next-token prediction on massive web-scale corpora.
    \item If scaling is predictable, then one can train many small models, fit empirical laws, and extrapolate to decide how to spend a much larger training budget.
    \item However, even if pretraining loss scales smoothly, the capabilities we care about, such as reasoning, safety, and tool use, may not behave as smoothly.
    \item The lecture therefore frames scaling both as a powerful engineering strategy and as an incomplete picture of progress.
\end{itemize}

\subsection{Key Concepts}
\textbf{Definition: Scaling Law}

A scaling law is an empirical power-law relationship between model performance and a resource such as parameter count, training tokens, or compute.

\begin{itemize}
    \item \textbf{Compute budget framing:} imagine having a fixed budget such as a specific cluster for a month. The question is which model/data/training configuration minimizes validation loss within that budget.
    \item \textbf{Validation loss and perplexity:} the lecture uses next-token prediction loss as the main axis of analysis because it is smooth, universal, and easy to measure.
    \item \textbf{Power law:} when plotted on log-log axes, many of these relationships become approximately linear.
    \item \textbf{Asymptotic thinking:} the important question is not how a model performs at a tiny scale, but how performance changes as scale increases.
    \item \textbf{Kaplan et al.\ scaling laws:} early large-scale evidence that loss decreases as a power law in model size, data size, and compute.
    \item \textbf{Chinchilla scaling:} a later refinement emphasizing compute-optimal tradeoffs between parameter count and number of training tokens.
    \item \textbf{Compute-optimal training:} the best use of a training budget may require balancing model size and data rather than maximizing only one of them.
    \item \textbf{Emergent behavior:} some capabilities appear to exhibit sudden qualitative jumps with scale, rather than only gradual smooth improvements.
    \item \textbf{In-context learning:} a model can adapt to a task from natural-language instructions and a few examples in the prompt, without gradient-based fine-tuning.
    \item \textbf{Chain-of-thought reasoning:} models can solve harder tasks by generating intermediate reasoning steps token by token.
    \item \textbf{Bitter Lesson perspective:} simple, general methods that scale with compute often outperform more hand-engineered, task-specific approaches.
\end{itemize}

\subsection{Core Models and Equations}
\textbf{Validation loss and perplexity.} Let \(L\) denote the average next-token prediction loss on a validation set. Then perplexity is
\[
\operatorname{PPL} = e^{L}.
\]
Perplexity can be interpreted as the average number of equally likely choices the model is uncertain between at each step.

\textbf{Generic power law.} The lecture frames a basic scaling relationship as
\[
L(X) \propto X^{-\alpha},
\]
where \(X\) is a resource such as number of parameters, training tokens, or compute, and \(0 < \alpha < 1\). Doubling the resource gives
\[
L(2X) = 2^{-\alpha}L(X),
\]
which means improvements are predictable but exhibit diminishing returns.

\textbf{Kaplan-style empirical scaling.} If data and compute are not bottlenecks, test loss scales approximately as a power law in the number of non-embedding parameters \(N\). Similarly, if model capacity is not the bottleneck, test loss scales approximately as a power law in the number of training tokens \(D\). Compute-optimal loss also follows a power law in total compute \(C\).

\textbf{Chinchilla parameter-data law.} A widely used formulation is
\[
L(N, D) = E + \frac{A}{N^{\alpha}} + \frac{B}{D^{\beta}},
\]
where:
\begin{itemize}
    \item \(E\) is the irreducible loss floor,
    \item \(N\) is model size,
    \item \(D\) is number of training tokens,
    \item \(A,B > 0\) are scale constants,
    \item \(\alpha,\beta \in (0,1)\) are power-law exponents.
\end{itemize}
The lecture emphasized that \(\alpha\) and \(\beta\) being less than 1 means diminishing returns.

\textbf{Approximate dense-transformer compute relation.} For rough planning, training compute can be modeled as
\[
C \approx 6 N D.
\]
This lets us connect model size and number of tokens under a fixed training budget.

\textbf{Compute-optimal choice.} Under a fixed \(C\), substitute
\[
D = \frac{C}{6N}
\]
into the Chinchilla loss model and minimize with respect to \(N\). This gives a recommended compute-optimal pair \((N,D)\) for the chosen budget.

\textbf{Why log-log plots help.} Taking logs of a power law
\[
L(X) = kX^{-\alpha}
\]
gives
\[
\log L(X) = \log k - \alpha \log X,
\]
which is linear in \(\log X\). This is why scaling-law papers visualize these relationships on log-log axes.

\subsection{Algorithms / Architectures}
\textbf{1. Naive strategies and why they fail}
\begin{itemize}
    \item Train many candidate models at equal compute and pick the best one.
    \item Problem: your final chosen model only used a fraction of your total compute budget.
    \item Tune hyperparameters at very small scale and copy them to large scale.
    \item Problem: good choices at small scale often do not transfer reliably to large scale.
\end{itemize}

\textbf{2. Scaling-law workflow}
\begin{itemize}
    \item Train many smaller models over a range of scales.
    \item Measure validation loss against compute, model size, and data size.
    \item Plot these relationships on log-log axes.
    \item Fit a power-law model.
    \item Extrapolate to the large training budget.
\end{itemize}

\textbf{3. Practical Chinchilla-style planning}
\begin{itemize}
    \item Choose several model sizes spanning at least an order of magnitude.
    \item For each size, try multiple learning-rate schedules and/or data budgets.
    \item Record validation loss at multiple compute points.
    \item Either fit the lower envelope of best-achievable loss versus compute, or directly fit
    \[
    L(N,D)=E+\frac{A}{N^\alpha}+\frac{B}{D^\beta}.
    \]
    \item Given a target compute budget \(C\), solve for the compute-optimal \(N\) and \(D\).
\end{itemize}

\textbf{4. Evaluating beyond perplexity}
\begin{itemize}
    \item Track validation loss and perplexity because they are smooth and universal.
    \item Also monitor downstream tasks such as general language understanding, coding, instruction following, reasoning, and live human usage.
    \item Ask whether loss reduction still correlates with the capabilities you care about.
\end{itemize}

\textbf{5. Interpreting emergence}
\begin{itemize}
    \item Measure capability as a function of scale.
    \item Look for thresholds where accuracy appears to jump sharply.
    \item Test whether the jump is robust or an artifact of the metric.
    \item Compare plain prompting, few-shot prompting, and chain-of-thought prompting.
\end{itemize}

\subsection{Important Intuitions from the Professor}
\begin{itemize}
    \item \textbf{Scaling is an engineering program, not a law of nature.} The professor used the analogy to Moore's Law: impressive exponential trends require continuous innovation and can eventually stall.
    \item \textbf{Loss is a very useful but incomplete proxy.} Smooth reductions in next-token prediction loss often correlate with downstream quality, but increasingly weakly for reasoning, safety, and other nuanced capabilities.
    \item \textbf{Kaplan's ``make the model bigger'' worldview was a productive illusion.} It was not literally correct, but it pushed the field toward valuable discoveries.
    \item \textbf{Chinchilla changed the optimization target.} The main message was not that scaling laws were wrong, but that compute-optimal allocation between parameters and tokens had been misestimated.
    \item \textbf{Data quality and architecture quality change the law you should fit.} One should not blindly copy coefficients from another group's paper if the dataset, optimizer, or model family is different.
    \item \textbf{Emergent behavior was genuinely surprising.} In-context learning and later chain-of-thought performance were not what many researchers expected from vanilla next-token prediction.
    \item \textbf{Emergence may partly reflect measurement artifacts.} A capability can look like a step function when the evaluation metric has thresholds or compounds multiple smaller skills.
    \item \textbf{Scaling alone is no longer the whole story.} Once a strong pretrained model exists, post-training, inference-time reasoning, tool use, and better inductive biases become more important levers.
\end{itemize}

\subsection{Examples}
\begin{itemize}
    \item \textbf{OpenAI 2018 AI-and-compute plot:} the lecture used this as an early public signal that training compute for frontier models had been growing rapidly.
    \item \textbf{Moore's Law analogy:} scaling in NLP was compared to decades of transistor growth in hardware.
    \item \textbf{Course-project framing:} the lecture asked what a student team would do with a fixed compute budget and why naive hyperparameter copying from small models is unreliable.
    \item \textbf{Kaplan et al.\ 2020:} showed approximate power-law relationships between test loss and parameter count, data size, and compute.
    \item \textbf{Chinchilla 2022:} argued that compute-optimal training requires scaling model size and training data more evenly than earlier views suggested.
    \item \textbf{Few-shot translation:} a sufficiently large autoregressive LM can see a few English-to-French examples in its context and continue with a new translation without gradient updates.
    \item \textbf{Chain-of-thought on arithmetic:} models can be prompted to write intermediate steps, and the capability appears to improve sharply with scale.
    \item \textbf{Grade School Math 8K (GSM8K) extrapolation caution:} naive extrapolation from pure scaling would have predicted absurdly large models to solve grade-school math, yet post-training and reasoning methods surpassed that expectation with much smaller systems.
\end{itemize}

\subsection{Common Pitfalls / Limitations}
\begin{itemize}
    \item \textbf{Overtrusting one fitted law:} coefficients depend on data quality, architecture, optimizer, and evaluation setup.
    \item \textbf{Optimizing only for training-time compute:} a deployment setting may prefer overtraining a smaller model if inference cost dominates over the model's lifetime.
    \item \textbf{Overfixating on perplexity:} it is not a sufficient proxy for reasoning, safety, or agentic usefulness.
    \item \textbf{Assuming emergence is always real:} some ``sudden jumps'' may be artifacts of thresholded or compounded evaluations.
    \item \textbf{Believing scale is all you need:} large models may be necessary for some capabilities, but they are far from sufficient.
    \item \textbf{Ignoring diminishing returns:} even if power laws continue, the extra compute required for small improvements may become unjustifiable.
    \item \textbf{Missing new axes of scaling:} post-training, tool use, and inference-time reasoning can matter as much as additional pretraining.
\end{itemize}

\subsection{Connections to Other Lectures}
\begin{itemize}
    \item This lecture builds directly on Lecture 6, which introduced pretraining and fine-tuning. Scaling laws answer how to choose the size of that pretraining run.
    \item It also connects to Lecture 5 on evaluation by questioning whether validation loss and perplexity are sufficient proxies for the capabilities we care about.
    \item The discussion of in-context learning and chain-of-thought connects back to Lecture 4's Transformer architecture, since these behaviors arise in large decoder-only Transformers trained with next-token prediction.
    \item The lecture sets up future lectures on post-training, reasoning, and tool use by arguing that pretraining loss reduction alone cannot explain current frontier-model progress.
\end{itemize}

\subsection{Exam-Relevant Summary}
\begin{itemize}
    \item Perplexity is
    \[
    \operatorname{PPL}=e^L,
    \]
    where \(L\) is validation loss.
    \item A basic scaling law has the form
    \[
    L(X)\propto X^{-\alpha},
    \]
    so doubling a resource gives predictable but diminishing improvements.
    \item Chinchilla-style scaling uses
    \[
    L(N,D)=E+\frac{A}{N^\alpha}+\frac{B}{D^\beta}.
    \]
    \item For dense Transformers, a rough compute model is
    \[
    C \approx 6ND.
    \]
    \item Kaplan et al.\ argued that bigger models were especially powerful; Chinchilla agreed on power laws but revised the compute-optimal allocation between parameters and data.
    \item Scaling-law experiments are done by training many smaller models, measuring loss, and extrapolating on log-log plots.
    \item You should refit scaling laws in your own setting if your data quality, architecture, or optimizer meaningfully differs.
    \item Loss and perplexity are smooth and useful, but they are increasingly weak proxies for capability, reasoning, and safety.
    \item Emergent behaviors discussed in this lecture include few-shot in-context learning and chain-of-thought reasoning.
    \item Some emergence claims may be artifacts of how capability is measured.
    \item Current progress goes beyond pure scaling: post-training, better inductive biases, reasoning strategies, and tool use are now major axes of improvement.
\end{itemize}

\end{document}
